% =====================================================================
% Appendix: Baseline implementation details and hyperparameters
% 본문 대응 문단 = \subsection{Baselines} 의 "Full implementation details and
% hyperparameters are provided in Appendix~\ref{app}."
%
% 근거 파일
%   EACL/UNATTENDED_2026-08-03_online_baselines.md            (D-01 ~ D-20)
%   outputs/experiments/2026-08-04_online_baselines_close/REPORT.md
%   outputs/experiments/2026-08-04_online_baselines_close/lookahead_audit.json
%   outputs/experiments/2026-08-03_narrative_online/table4_6docs.json
%   scripts/eacl_phase0/make_main_table.py                    (표 행 구성·명명 규칙)
% =====================================================================

\section{Baseline Implementation Details}
\label{app:baselines}

This appendix documents every row of Table~\ref{tab:main}: what code and weights
were run, which hyperparameters were used, what had to be changed to satisfy the
no-future-context setting, and which variants were measured but excluded.
All numbers referenced here were produced by the runners named in each
paragraph; predictions for every system are dumped in a common format
(one binary vector of length $n-1$ per document, indexed by transition) and
scored by a single scorer, so that no row is advantaged by a different
evaluation path.

% ---------------------------------------------------------------------
\subsection{Shared evaluation protocol}
\label{app:baselines-protocol}

\paragraph{Evaluation data.}
All systems are evaluated on the same six narratives, three development and
three held-out (Table~\ref{tab:app-eval-data}). The development narratives are
the spoken accounts used by \citet{michelmann2025large}; the held-out narratives
are three excerpts from \emph{Born a Crime} annotated by 20 human listeners,
whose consensus boundaries form the gold standard. Boundary density varies by a
factor of five across documents ($0.036$--$0.190$), which is why every claim in
the main table is read against a no-boundary floor computed per document set.

\begin{table}[t]
\centering\small
\begin{tabular}{@{}llrrr@{}}
\toprule
Split & Document & Sent. & Trans. & Gold \\
\midrule
\multirow{3}{*}{dev-3}
 & Pieman        &  96 &  95 & 18 \\
 & Tunnel        & 504 & 503 & 18 \\
 & Monkey        & 408 & 407 & 30 \\
\addlinespace
\multirow{3}{*}{test-3}
 & Run           & 150 & 149 & 14 \\
 & GoHitler      & 113 & 112 & 13 \\
 & MyMothersLife & 172 & 171 & 16 \\
\midrule
 & Total         & 1{,}443 & 1{,}437 & 109 \\
\bottomrule
\end{tabular}
\caption{Evaluation data. ``Trans.'' is the number of scored transitions
($=$ sentences $-1$); ``Gold'' is the number of human boundaries.}
\label{tab:app-eval-data}
\end{table}

\paragraph{Boundary index convention.}
Transition $j$ is the position between sentence $j$ and sentence $j+1$
($0$-indexed). Streaming segmenters that report ``the utterance that opens a new
segment'' (index $b$, $1$-based) are mapped to transition $b-1$. Following the
convention of \citet{jiang2023superdialseg}, the last transition of a document
is always labelled $0$ for both gold and prediction.

\paragraph{Deployable scoring.}
Every system decides for itself how many boundaries a document contains; the
gold boundary count is never supplied to any method. Where a method's published
read-out requires an externally supplied count (top-$k$), we replace it and say
so explicitly (\S\ref{app:baselines-dts}). Thresholds and adaptation rules are
frozen on dev-3 and never re-selected on test-3.

\paragraph{Metrics.}
Score, BOR and Lift are as defined in \S\ref{sec:metrics}. Score is computed per
document and averaged; the per-split aggregates quoted in this appendix are the
mean over the three documents of that split, and the ``Avg'' column of
Table~\ref{tab:main} is the mean over all six.
AUC is not reported for baselines: their prediction dumps are binary, so no
continuous ranking signal is available.

% ---------------------------------------------------------------------
\subsection{The online criterion and how it was audited}
\label{app:baselines-online}

\paragraph{Definition.}
A system is \emph{online} (zero-lookahead) if the decision for transition $j$
uses only sentences $u_{\le j+1}$ --- that is, the decision is final at the
moment the sentence that would open the new segment arrives. A system is
\emph{bounded-lookahead with delay $d$} if the decision for transition $j$ is
emitted only after $u_{j+1+d}$ has been observed.

\paragraph{Measured, not assumed.}
Push-based streaming segmenters trivially pass a future-perturbation test,
because future sentences have not been pushed yet; such a test is therefore no
evidence at all. What actually varies is \emph{latency}. We instrumented every
push-based segmenter: if \texttt{push}$(u_t)$ emits boundary $b$, the decision
for transition $b-1$ used $u_b \dots u_t$, so its delay is $t-b$; we take the
maximum over all six documents. Two implementations turned out to advertise the
wrong grade, in both cases optimistically:

\begin{itemize}\itemsep2pt
\item A class named \texttt{GreedySegOnlineDelay2} (suggesting $d{=}2$) measured
      $d{=}7$, because its finalisation routine accumulates candidates and
      commits them in a batch.
\item A graph-based segmenter configured with a window of $10$ measured $d{=}6$,
      because it does not emit until its window is full.
\item The streaming TextTiling variant that our earlier tables listed simply as
      ``TextTiling'' measured $d{=}22$. Its advantage over the zero-lookahead
      version ($+0.057$ Score on dev-3) is entirely an advantage of seeing 22
      sentences of future context.
\end{itemize}

\paragraph{Consequence for the main table.}
Because a ``bounded'' label without a number is uninformative --- readers
imagine $d{=}2$ when the truth may be $d{=}22$ --- Table~\ref{tab:main} contains
\emph{only} rows whose measured maximum delay is $0$. Where the published method
is structurally bounded, we re-ran it in a zero-lookahead configuration and
marked it as a deviation from the original (\S\ref{app:baselines-naming}); the
bounded originals are reported in \S\ref{app:baselines-excluded} so that the
cost of removing lookahead is visible.

% ---------------------------------------------------------------------
\subsection{Naming convention in Table~\ref{tab:main}}
\label{app:baselines-naming}

\begin{itemize}\itemsep2pt
\item \textbf{\texttt{X}} --- the published method run as published
      (offline). Only TextTiling appears in this form.
\item \textbf{\texttt{X-style Online}} --- the published method was offline or
      symmetric-context, and we re-ran its scoring function under the
      zero-lookahead read-out described below.
\item \textbf{plain name} (EM-LLM, BayesSurprise) --- the published method is
      already causal; nothing was changed, so no suffix is added.
\item $\dagger$ --- the adaptation changes the \emph{method definition}, not
      only the inference schedule (Embedding Delta's symmetric window,
      DialSTART's top-$k$ read-out). These rows are not comparable to the
      corresponding published numbers.
\item $\ddagger$ --- not a published model: our own re-implementation trained on
      our anchor labels (\S\ref{app:baselines-ssc}).
\end{itemize}

% ---------------------------------------------------------------------
\subsection{Unsupervised narrative event segmentation}
\label{app:baselines-nes-unsup}

\paragraph{Embedding Delta \citep{schneider2021detecting} ---
\texttt{EmbeddingDelta-style Online}$^{\dagger}$.}
The original pipeline embeds tokens with spaCy vectors, fits a $k$-means
codebook over the \emph{whole document}, represents each sentence as a cluster
histogram, computes a normalised cosine distance between a left and a right
window of 25 sentences each, smooths the signal with a moving average of 10, and
takes local maxima with \texttt{argrelmax(order=20)} as boundaries. Three of
these steps use the future and all three were changed:

\begin{table}[t]
\centering\small
\begin{tabular}{@{}p{0.30\columnwidth}p{0.28\columnwidth}p{0.30\columnwidth}@{}}
\toprule
Original & Future used & Online replacement \\
\midrule
$k$-means over all document tokens & whole document &
\texttt{MiniBatchKMeans.partial\_fit} on tokens seen so far \\
right window of $w$ sentences & $w$ ahead &
right window of $1$ sentence (\texttt{-{}-right-w 1}) \\
\texttt{argrelmax(order=20)} & $20$ ahead &
running $\mu + c\sigma$ over past scores \\
\bottomrule
\end{tabular}
\caption{Making Embedding Delta zero-lookahead.}
\label{tab:app-embdelta}
\end{table}

Hyperparameters follow the original where they are causal: $k=1000$ clusters,
left window $w=25$ sentences, smoothing window $10$, threshold constant
$c=0.5$ (the original's ``$\mathrm{mean}+\mathrm{std}/2$'' rule), warm-up $5$
scored transitions, spaCy \texttt{en\_core\_web\_lg} vectors,
\texttt{MiniBatchKMeans}$(n\_init{=}3, \text{batch}{=}100,
\texttt{random\_state}{=}0)$. The \emph{left} window is deliberately kept at
$w=25$: shrinking both windows would confound ``worse because zero-lookahead''
with ``worse because the window is narrow''.

The right-window reduction is the reason for the $\dagger$: symmetric context is
part of the original method definition. Two consequences are worth recording.
First, removing lookahead does \emph{not} cost anything here --- the
zero-lookahead version is $+0.014$ Score better on dev-3 and $-0.006$ worse on
test-3 than the $d{=}25$ original, in contrast to TextTiling, where the same
move costs $-0.057$ on dev-3. Lexical-cohesion signals profit from future
context; codebook-histogram signals apparently do not. Second, the symmetric
version needs 50 sentences before it can score anything and therefore emits
\emph{zero} predictions on the 96-sentence Pieman; the zero-lookahead version
needs 26 and emits 7.

\emph{Reproducibility note.} The original code path had no seed on
\texttt{MiniBatchKMeans}, so the codebook --- and the score --- differed between
identical runs (dev-3 Score $0.2811 / 0.2809 / 0.2810$). We added
\texttt{random\_state} and verified that two runs of the same configuration now
agree exactly. Numbers reported before this fix should be read as
$\pm 0.004$.

% ---------------------------------------------------------------------
\subsection{Supervised narrative event segmentation}
\label{app:baselines-ssc}

\paragraph{Why no published model can be run.}
The three supervised scene-segmentation models in the narrative literature ---
sequential sentence classification with half-stride inference
\citep{zehe2025assessing}, the sentence-pair classifier of
\citet{zehe2021detecting}, and the transfer approach of
\citet{kurfali2021breaking} --- are German-only: they are built on German
encoders and trained on the German Scene Dataset. Our evaluation narratives are
English, so even a working checkpoint would have no zero-shot path.
Independently, the released checkpoint for \citet{kurfali2021breaking} is no
longer retrievable (the distribution link returns HTTP 404), the training corpus
of \citet{zehe2025assessing} is copyrighted and not redistributable, and
\citet{barth2021participation} requires a licensed lexical resource. We did not
contact authors or request licences.

\paragraph{\texttt{SSC-style Online (our anchors)}$^{\ddagger}$.}
We therefore keep the \emph{architecture} and replace the training data: a
sequential sentence classifier that reads a sliding context window of sentences
and classifies each sentence as segment-final or not. Backbone
\texttt{roberta-base} (the same backbone as the supervised DTS baseline, so the
comparison is not confounded by encoder capacity), token-classification head
with two labels. Training data are our anchor labels --- the same distant
supervision that M4/M5/M8 receive, never human boundary annotations. Training
configuration: AdamW, learning rate $1\times10^{-5}$, weight decay $1\times
10^{-3}$, gradient clipping $1.0$, batch $8$, at most 20 epochs with early
stopping on validation loss (patience 3) and best-epoch restore, bf16 autocast,
seed 42; sliding window 20 sentences, at most 25 word-pieces per sentence, at
most 512 word-pieces per window; documents split 90/10 into train and
validation.

Inference is the zero-lookahead path: the boundary between sentences $t-1$ and
$t$ is decided once, when $t$ arrives, from the causal window
$u_{\max(0,t-W+1)\dots t}$, classifying the position of sentence $t-1$. In
particular \emph{Half Stride} --- the inference procedure that
\citet{zehe2025assessing} introduce precisely because one-sided context is
insufficient --- is not used, since it averages predictions over windows
extending into the future. The gap between our online row and their reported
offline performance should be read in that light; we cite their argument but do
not report their numbers, since they are not comparable.

This row is an ablation of our own, not a baseline: it shares our training
signal and differs only in architecture. It is marked $\ddagger$ and described
as ``our re-implementation, trained on our anchor labels'' wherever it appears.

% ---------------------------------------------------------------------
\subsection{LLM-signal narrative event segmentation}
\label{app:baselines-llm}

\paragraph{EM-LLM \citep{fountas2025human}.}
We implement the surprise-threshold (SM) rule exactly as in the reference
implementation: a transition is a boundary when the sentence-level surprise
exceeds the mean plus $\gamma$ standard deviations of the surprise observed in
the preceding window,
$s_t > \mu(\text{win}) + \gamma\,\sigma(\text{win})$, with the window defined
over the previous $n_{\text{local}}$ \emph{tokens} (not sentences) of the same
document. Defaults are taken from the released configuration:
$\gamma = 1.1$, $n_{\text{local}} = 4096$ tokens, no per-document tuning; a
warm-up of 3 in-window sentences before the first decision.

Two scoping decisions matter. (i) The graph-refinement stage (SM$+$), which
re-places boundaries using modularity/conductance over the similarity graph of
the \emph{following} segment, is excluded: it requires future context and would
break the online claim. (ii) Sentence-level surprise is computed with a causal
language model, and the tokenizer used to count $n_{\text{local}}$ must be that
same model's tokenizer --- an early version of our runner counted a Qwen window
with a Llama tokenizer. Correcting this did not change the scores (window
coverage was already $\approx 1.0$), but the runners were fixed.

The main-table row uses \texttt{Qwen2.5-7B} as the surprise model. Because
EM-LLM's own experiments use Llama, we recomputed the entire pipeline with two
Llama models to rule out a model-choice artefact
(Table~\ref{tab:app-emllm-lm}); all three fall below the no-boundary floor on
both splits, and the Llama models --- including the exact model used in the
original work --- are worse than Qwen.

\begin{table}[t]
\centering\small
\begin{tabular}{@{}lrr@{}}
\toprule
Surprise model & dev-3 & test-3 \\
\midrule
\texttt{Qwen2.5-7B} (main table)     & 0.2699 & 0.2784 \\
\texttt{Llama-3.1-8B}                & 0.2511 & 0.2425 \\
\texttt{Llama-3.1-8B-Instruct}       & 0.2230 & 0.2371 \\
\midrule
\emph{no-boundary floor}             & \emph{0.2790} & \emph{0.3092} \\
\bottomrule
\end{tabular}
\caption{EM-LLM Score under three surprise models, all other settings fixed
($\gamma{=}1.1$, $n_{\text{local}}{=}4096$, SM only). The segmentation code is
byte-identical across rows; only the language model that produces the surprise
sequence changes.}
\label{tab:app-emllm-lm}
\end{table}

\emph{Implementation notes.} The encoder must be loaded on a single GPU:
sharding it with \texttt{device\_map="auto"} splits layers across devices and
the per-layer representation stack then mixes tensors from different devices.
When GPU memory is shared with other jobs, token negative log-likelihoods are
computed in row chunks; \texttt{log\_softmax} is row-independent, and we
verified $\max|\Delta s| = 0$ between the chunked and unchunked paths on one
document before using it.

\paragraph{BayesSurprise \citep{kumar2023bayesian}.}
Word-level Bayesian surprise is the KL divergence between consecutive
next-word distributions, $S_i = \mathrm{KL}(P(\cdot\mid w_{\le i}) \,\|\,
P(\cdot\mid w_{\le i-1}))$, taken at the last word-piece of each word, computed
with GPT-2-XL under a moving 1024-token context (stride 512), in float32
(the KL is precision-sensitive). Surprise is transient-normalised following
Reynolds: $\mu_1 = \overline{S}$, $\mu_t = \mu_{t-1} + (S_t-\mu_{t-1})\delta$,
$\mathrm{Transient}_t = S_t/\mu_{t-1}$, with drift $\delta = 0.10$.

The read-out is \texttt{newsent\_max}: the score of transition $t$ is the
maximum transient surprise over all words of sentence $t+1$. The threshold is a
running $\mu + c\sigma$ over the scores observed so far in the same document,
with $c = 0.75$ frozen on Pieman and never re-tuned, and a warm-up of 3
transitions. The method is causal by construction: normalisation uses only a
past moving average, and the threshold uses only scores already observed.

\emph{Read-out choice, recorded.} An onset-max read-out (maximum over the first
$N$ words of the new sentence) was implemented first and is \emph{not} used: it
is anti-correlated with boundaries in spoken narrative (AUC $0.35$), because
event openings are formulaic (``And so\ldots'') and therefore low-surprise.
Using it would have cost $-0.048$ dev-3 / $-0.019$ test-3.

\paragraph{EventRecall \citep{panela2025event} ---
\texttt{EventRecall-style Online}.}
We port only the segmentation stage of that pipeline; the recall-scoring
contribution is out of scope and this must be stated whenever the method is
cited. The published procedure sends the entire narrative to the LLM in a single
call and parses the returned event blocks --- an offline procedure. The online
version asks, for each arriving sentence, a single binary question: given the
narrative up to and including the current sentence, and given the segmentation
the model has \emph{already committed to}, does this sentence begin a new event
or continue the current one? The event criterion, the wording of the
segmentation goal and the output format are taken verbatim from the original
prompt; only the question is converted from ``partition this text'' to ``NEW or
SAME''.

Settings: \texttt{gpt-4.1} via OpenRouter, temperature $0$, at most 4 completion
tokens, full past context (no truncation), 10 independent calls per transition
with majority vote (matching the repetition-and-consensus aggregation of the
offline procedure), stateful prompt (previously committed segment ids are shown,
earlier decisions are final). Decisions within a document are therefore
sequential; no future sentence is ever shown. Responses are cached, so the
reported row is a single deterministic replay of one API run.

\emph{Not comparable to the published numbers.} The original work reports a
Hamming-distance-based segmentation figure that is dominated by boundary
sparsity on our gold and that aggregates gold differently; there is no published
number this row can be matched against.

% ---------------------------------------------------------------------
\subsection{Dialogue topic segmentation}
\label{app:baselines-dts}

\paragraph{TextTiling \citep{hearst1997text} --- two rows.}
The \texttt{TextTiling} row is the unmodified NLTK implementation
(\texttt{TextTilingTokenizer}, pseudo-sentence size $w=10$, block size $k=6$),
applied to the whole document at once; sentences are joined by blank lines and
the tile boundaries are mapped back to transitions. It is offline and is
included as the reference point for what removing future context costs.

The \texttt{TextTiling-style Online} row is the same scoring function under
prefix re-computation: the decision for transition $j$ re-runs the tokenizer on
$u_{\le j+1}$ and reads off position $j$. This is deliberately \emph{not} a
streaming approximation --- the only variable that changes relative to the
offline row is the availability of future context, so the difference is
attributable to that alone. The cost is $O(n^2)$ tokenizer calls per document;
on the 503-transition Tunnel this is minutes rather than the milliseconds the
streaming variant takes, which is itself a finding: deploying a zero-lookahead
lexical-cohesion segmenter requires an incremental implementation, not
re-computation.

\paragraph{DialSTART \citep{gao2023unsupervised} ---
\texttt{DialSTART-style Online}$^{\dagger}$.}
Weights are the authors' released \texttt{model\_711} checkpoint; the coherence
scoring model, the tokenisation, the \texttt{window\_size} parameter and all
weights are used unchanged. Three modifications were required and all three are
marked by the $\dagger$:

\begin{enumerate}\itemsep2pt
\item \textbf{Read-out.} The released inference selects boundaries by
      \texttt{argsort(depth)[-pick\_num:]} with $\texttt{pick\_num}=4$ --- a
      \emph{fixed} number of boundaries per document, independent of document
      length. This is not gold leakage, but it is incompatible with deployable
      scoring, and on documents with 95--503 transitions and 13--30 gold
      boundaries it is a severe under-prediction. We replace it with a running
      $\mu + c\sigma$ threshold.
\item \textbf{Depth computation.} The original right-peak search scans to the
      end of the document. We bound it (\texttt{depth\_span}$=0$ in the
      main-table configuration).
\item \textbf{Scoring window.} The original scores transition $i$ from
      $\texttt{text}[i{+}1], \texttt{text}[i{+}2]$. We reduce the right side to
      the arriving sentence only (\texttt{-{}-right-window 1}), keeping the left
      side at the published \texttt{window\_size}$=2$.
\end{enumerate}

Together these give a measured maximum delay of $0$. The threshold constant is
$c=0.5$, matching the sign convention used by the other unsupervised rows
(streaming TextTiling uses $c=0.5$; Embedding Delta uses the original's
$\mathrm{mean}+\mathrm{std}/2$). \textbf{$c$ is not tuned.} An earlier
configuration froze $c$ by sweeping on dev-3, which reached a much higher score
($0.3586$ dev-3 / $0.3462$ test-3) --- but a baseline whose hyperparameter is
fitted to our data is no longer that baseline, it is a variant of our own
method, so the published constant is used instead. Untuned, the method
over-segments badly on narrative (290 predictions against 66 gold boundaries on
dev-3, BOR $3.62$).

\emph{Not comparable to the published numbers.} The original reports Pk/WD on
dialogue corpora with the top-$k$ read-out; this row measures what happens when
a dialogue-trained segmenter is transferred zero-shot to narrative.

\paragraph{Supervised RoBERTa \citep{jiang2023superdialseg} ---
\texttt{RoBERTa-style Online}.}
A \texttt{roberta-base} token classifier trained on the SuperDialseg training
split: AdamW, learning rate $1\times10^{-5}$, batch 8, weight decay
$1\times10^{-3}$, gradient clipping $1.0$, 20 epochs with early stopping on
validation Score (patience $=$ epochs$/2$, best checkpoint retained), no LR
schedule, sliding window of 20 utterances with stride 1, at most 25 word-pieces
per utterance, at most 512 per window, seed 42. Weights are unchanged from that
training run; only inference differs. Offline inference decides utterance $g$
by averaging the logits of every stride-1 window containing $g$, including
windows reaching up to 19 utterances into the future. Online inference decides
the boundary $(t-1,t)$ once, when $t$ arrives, from the single causal window
$u_{\max(0,t-W+1)\dots t}$, classifying the second-to-last utterance --- a
position that carried a live label during training. Lookahead is zero by
construction (the slice upper bound is $t+1$), one window per turn, no revision.
On narrative, sentences are treated as utterances.

\paragraph{CSM \citep{xing2021improving} --- \texttt{CSM-style Online}.}
Coherence scores come from the original coherence network (BERT backbone plus
MLP head) applied to adjacent sentence pairs, using our own re-trained
checkpoint. Depth scores use the TextTiling-style left/right peak construction,
and the threshold is the original $\mathrm{mean} + \alpha\,\mathrm{std}$ with
$\alpha=1.0$, maintained online by a Welford running estimator, plus the
original minimum-gap suppression ($\texttt{min\_gap}=2$). The published
inference computes both peaks over the full score sequence; the streaming
variant we had used previously bounded the right peak to $d=2$. For the main
table we set $d=0$, i.e.\ the right peak is dropped entirely and the depth at a
transition is computed from the left peak alone. This is not a loss:
$d{=}0$ scores slightly \emph{above} $d{=}2$ (Avg $0.312$ vs $0.308$).

\emph{Checkpoint caveat.} The authors' released checkpoint is not distributed
(the reproduction script references a file that has no download location) and
our earlier local copy was lost. The checkpoint used here is our own
re-training on DailyDialog. That re-training is known to differ from the
released model: the upstream pseudo-sample generator has no seed (sample counts
varied by $\sim0.4\%$ across runs) and we fixed three upstream bugs, notably a
validation routine that never restored training mode, leaving dropout disabled
for 24k steps. Against the authors' own benchmarks this checkpoint reproduces
two of three datasets (TIAGE $99.9\%$, SuperDialseg $110.1\%$ of the published
score) but reaches only $80.4\%$ on DialSeg711 at its best epoch. The shortfall
is attributable to checkpoint availability, not to the method, and must be
stated wherever this row is discussed.

% ---------------------------------------------------------------------
\subsection{Controls}
\label{app:baselines-controls}

All three controls are online by definition: they use only position, a
pseudo-random draw, or nothing at all.

\begin{itemize}\itemsep3pt
\item \textbf{Fixed interval (density-matched).} $k$ boundaries placed at
      equal spacing over $n$ transitions, deterministic.
\item \textbf{Random (density-matched).} $k$ of the $n$ transitions chosen
      uniformly without replacement. The draw is keyed by transition index, so
      re-asking about an earlier transition returns the same answer --- required
      for the causality audit to be meaningful. Twenty replicates were generated
      (seed \texttt{20260803}); the table reports the first, and the dev-3 mean
      over replicates is $0.2763 \pm 0.0302$, so the single-replicate row should
      not be read as precise to the third decimal.
\item \textbf{No boundary.} Predicts nothing. Because $P_k = \mathrm{WD}$ when
      no boundary is predicted, Score collapses to $0.5(1-P_k)$; this is the
      degenerate floor against which every row must be read.
\end{itemize}

\emph{What the density is matched to.} $k$ is set per document to the number of
boundaries \emph{our own model} predicts on that document, not to the gold
count --- matching gold would leak the oracle count into the control. On dev-3
the reference is the M5 prediction dump; on test-3, where no M5 dump exists, the
reference is the M5-adaptive dump. The controls therefore isolate the
contribution of \emph{where} boundaries are placed, holding \emph{how many}
fixed at our own operating point.

\emph{Note.} These are not the same as the similarly named controls in the
SuperDialseg toolkit, which draw the boundary \emph{count} at random and hence
do not control density.

% ---------------------------------------------------------------------
\subsection{Variants measured but excluded from the main table}
\label{app:baselines-excluded}

Table~\ref{tab:app-excluded} lists configurations that were run and scored but
are not in Table~\ref{tab:main}, together with the reason. They are reported
because two of them are \emph{better} than the row we kept, and omitting them
silently would misrepresent the comparison.

\begin{table}[t]
\centering\small
\begin{tabular}{@{}llrl@{}}
\toprule
Method & Config & Avg & Excluded because \\
\midrule
TextTiling      & bounded $d{=}22$      & 0.360 & uses 22 sentences of future \\
Embedding Delta & symmetric $d{=}25$    & 0.315 & uses 25 sentences of future \\
CSM             & bounded $d{=}2$       & 0.308 & uses future \\
DialSTART       & $d{=}2$, $c$ tuned    & 0.320 & uses future; $c$ fitted to dev-3 \\
EM-LLM          & Llama-3.1-8B          & 0.247 & LM-choice control, not the main config \\
EM-LLM          & Llama-3.1-8B-Instruct & 0.230 & LM-choice control \\
GraphSeg        & bounded $d{=}6$       & --    & no zero-lookahead form exists \\
GreedySeg       & bounded $d{=}7$       & --    & no zero-lookahead form exists \\
\bottomrule
\end{tabular}
\caption{Configurations excluded from Table~\ref{tab:main}, with their six-document
mean Score where applicable. GraphSeg and GreedySeg are dropped entirely: their
window-graph construction and their batched candidate finalisation are the
algorithms themselves, so a zero-lookahead version would no longer be the
published method.}
\label{tab:app-excluded}
\end{table}

Additional LLM-prompting variants (GPT-5 and \texttt{gpt-oss-120b} versions of
the online EventRecall prompt, a memory-system topic-shift prompt, and a
dialogue-oriented segmentation prompt) were also run; they are omitted to keep
one representative row per method family.

% ---------------------------------------------------------------------
\subsection{Methods we could not run, and why}
\label{app:baselines-unreproducible}

\begin{table}[t]
\centering\small
\begin{tabular}{@{}p{0.34\columnwidth}p{0.58\columnwidth}@{}}
\toprule
Method & Blocking reason \\
\midrule
SSC / half-stride \citep{zehe2025assessing} &
German-only encoder and corpus; training corpus is copyrighted and not
redistributable \\
Sentence-pair classifier \citep{zehe2021detecting} &
same corpus; no released checkpoint \\
Transfer SSC \citep{kurfali2021breaking} &
released checkpoint returns HTTP 404; no mirror \\
Lexical-chain segmentation \citep{barth2021participation} &
requires a licensed German lexical resource \\
Direct LLM prompting \citep{michelmann2025large} &
the reported model has been retired; no comparable number remains \\
\bottomrule
\end{tabular}
\caption{Published methods that could not be reproduced. We did not contact
authors or request licences; the reasons are recorded instead.}
\label{tab:app-unrepro}
\end{table}

For \citet{zehe2025assessing} we cite the \emph{argument} that motivated their
half-stride inference --- that one-sided context is insufficient for reliable
boundary decisions --- without reporting their numbers, and we substitute the
re-implementation of \S\ref{app:baselines-ssc}, clearly marked as ours.

% ---------------------------------------------------------------------
\subsection{Reproducibility and known limitations}
\label{app:baselines-limits}

\begin{itemize}\itemsep3pt
\item \textbf{Single harness.} Every number in Table~\ref{tab:main} is produced
      by one scorer over one prediction-dump format. Numbers from the original
      authors' harnesses are never mixed into that table: for the same method
      and the same checkpoint, moving between harnesses moved TIAGE Score from
      $0.3081$ to $0.3628$, which is larger than most of the differences we
      discuss.
\item \textbf{Our rows and the baseline rows come from different pipelines.}
      The M5/M4/M8 rows are means over four seeds and are not recomputed from
      the shared prediction dumps, whereas every baseline row is a single run
      recomputed from them. Per-seed dumps should be written to the shared
      location and the whole table recomputed by one runner.
\item \textbf{Anchor-corpus version.} The SSC re-implementation was trained on
      the anchor label index containing 10{,}598 documents and 2{,}777{,}561
      transitions (boundary rate $0.0657$ at the frozen threshold of three of
      five dimensions changing). The training-setup description in
      \S\ref{sec:experimental-design} quotes a later build of the anchor corpus.
      The claim that this row shares our training signal is therefore accurate
      in kind but not in corpus version.
\item \textbf{Split protocol.} The SSC re-implementation splits train/validation
      at the document level; our own models split by author. This makes its
      validation split easier than theirs.
\item \textbf{Three development documents.} Every frozen operating point
      ($c$, thresholds, read-out choices) was selected on three documents. For
      DialSTART the dev-3 sweep was nearly flat over $c\in[1.0,2.0]$
      ($0.3254$--$0.3295$), so such a selection would have been unstable; this
      is one reason the published constant is used instead.
\item \textbf{Seed sensitivity.} Embedding Delta is seed-sensitive through its
      codebook (spread $\approx 0.004$ Score). A multi-seed mean would be more
      appropriate than the single seed reported.
\item \textbf{No AUC for baselines.} Prediction dumps are binary. Dumping the
      continuous score alongside the binary decision would make ranking-based
      diagnostics available for every row.
\end{itemize}
